\documentclass[11pt,a4paper]{article}
\usepackage[utf8]{inputenc}
\usepackage{amsmath,amssymb,amsthm}
\usepackage{listings}
\usepackage{xcolor}
\usepackage[margin=1in]{geometry}
\usepackage{hyperref}

\newtheorem{theorem}{Theorem}
\newtheorem{lemma}[theorem]{Lemma}

\lstset{
  language=C++,
  basicstyle=\ttfamily\small,
  keywordstyle=\color{blue}\bfseries,
  commentstyle=\color{green!60!black},
  stringstyle=\color{red},
  numbers=left,
  numberstyle=\tiny\color{gray},
  breaklines=true,
  frame=single,
  tabsize=4,
  showstringspaces=false
}

\title{IOI 2002 -- Batch Scheduling}
\author{Solution}
\date{}

\begin{document}
\maketitle

\section{Problem Statement Summary}
There are $N$ jobs ($N \le 10{,}000$) to be processed on a single
machine in the fixed order $1, 2, \ldots, N$. Jobs may be grouped into
consecutive batches. Each batch incurs a setup time of $S$ before
processing begins. Job $i$ has processing time $t_i$ and cost weight
$f_i$. All jobs in a batch finish simultaneously at the batch's
completion time. The cost of job $i$ is $f_i \cdot C_i$, where $C_i$
is the completion time of job $i$'s batch. Minimize
$\sum_{i=1}^{N} f_i \cdot C_i$.

\section{Solution: DP with Convex Hull Trick}

\subsection{DP Formulation}
Define prefix sums $T[j] = \sum_{i=1}^{j} t_i$ and
$F[j] = \sum_{i=1}^{j} f_i$, with $T[0] = F[0] = 0$.

Let $dp[j]$ be the minimum total cost for scheduling jobs
$1, \ldots, j$, where job $j$ is the last job in its batch.

If the last batch contains jobs $k{+}1, \ldots, j$, then:
\begin{itemize}
  \item The batch finishes at time
        $C = S + T[j]$ plus the total setup and processing time of all
        previous batches.
  \item The setup time $S$ for this batch delays \emph{all} jobs
        $k{+}1, \ldots, N$.
\end{itemize}

A standard reformulation accounts for the delay penalty directly:
\[
  dp[j] = \min_{0 \le k < j}
  \Bigl\{ dp[k] + \bigl(S + T[j] - T[k]\bigr)
  \cdot \bigl(F[N] - F[k]\bigr) \Bigr\}.
\]
Base case: $dp[0] = 0$. Answer: $dp[N]$.

\begin{lemma}[Correctness of the reformulation]
Expanding the recurrence telescopically, $dp[N]$ equals the total
weighted completion time $\sum_{i=1}^{N} f_i \cdot C_i$.
The key idea is that placing setup time $S$ before the batch ending at
$j$ delays all remaining jobs $k{+}1, \ldots, N$ by exactly
$S + (T[j] - T[k])$, contributing cost proportional to
$F[N] - F[k]$.
\end{lemma}

\subsection{Convex Hull Trick}
Expand the transition:
\begin{align*}
  dp[j] &= \min_{k} \Bigl\{
    dp[k] + (S + T[j] - T[k]) \cdot (F[N] - F[k])
  \Bigr\} \\
  &= \min_{k} \Bigl\{
    \underbrace{(F[N] - F[k])}_{\text{slope}} \cdot T[j]
    + \underbrace{dp[k] + (S - T[k])(F[N] - F[k])}_{\text{intercept}}
  \Bigr\}.
\end{align*}
Define line $\ell_k: y = m_k \cdot x + b_k$ where:
\[
  m_k = F[N] - F[k], \qquad
  b_k = dp[k] + (S - T[k]) \cdot (F[N] - F[k]).
\]
Since $F[k]$ is non-decreasing, the slopes $m_k$ are non-increasing.
The query points $x = T[j]$ are non-decreasing.

We maintain a convex hull (lower envelope) of lines in a deque. Lines
are added at the back (decreasing slope); queries advance from the
front (increasing $x$). Each line is inserted and removed at most once,
giving amortized $O(1)$ per transition.

\section{C++ Implementation}

\begin{lstlisting}
#include <bits/stdc++.h>
using namespace std;

int main() {
    ios::sync_with_stdio(false);
    cin.tie(nullptr);

    int N;
    long long S;
    cin >> N >> S;

    vector<long long> t(N + 1), f(N + 1);
    vector<long long> T(N + 1, 0), F(N + 1, 0);

    for (int i = 1; i <= N; i++) {
        cin >> t[i] >> f[i];
        T[i] = T[i - 1] + t[i];
        F[i] = F[i - 1] + f[i];
    }

    vector<long long> dp(N + 1, 0);

    // Convex Hull Trick: minimise over lines y = m*x + b
    // with decreasing slopes and increasing query points.
    struct Line {
        long long m, b;
        long long eval(long long x) const { return m * x + b; }
    };

    deque<Line> hull;

    auto bad = [](const Line& l1, const Line& l2,
                  const Line& l3) -> bool {
        return (__int128)(l3.b - l1.b) * (l1.m - l2.m)
            <= (__int128)(l2.b - l1.b) * (l1.m - l3.m);
    };

    auto addLine = [&](long long slope, long long intercept) {
        Line ln{slope, intercept};
        while (hull.size() >= 2 &&
               bad(hull[hull.size() - 2], hull.back(), ln))
            hull.pop_back();
        hull.push_back(ln);
    };

    auto query = [&](long long x) -> long long {
        while (hull.size() >= 2 &&
               hull[0].eval(x) >= hull[1].eval(x))
            hull.pop_front();
        return hull.front().eval(x);
    };

    // Add line for k = 0
    addLine(F[N], S * F[N]); // m=F[N]-F[0]=F[N], b=dp[0]+(S-0)*F[N]

    for (int j = 1; j <= N; j++) {
        dp[j] = query(T[j]);

        long long slope = F[N] - F[j];
        long long intercept = dp[j] + (S - T[j]) * (F[N] - F[j]);
        addLine(slope, intercept);
    }

    cout << dp[N] << "\n";

    return 0;
}
\end{lstlisting}

\section{Complexity Analysis}
\begin{itemize}
  \item \textbf{Time:} $O(N)$. Each line is inserted into and removed
        from the deque at most once.
  \item \textbf{Space:} $O(N)$.
\end{itemize}

\end{document}
